\section*{5.5 计算梯度的常用恒等式}
\addcontentsline{toc}{section}{5.5 计算梯度的常用恒等式}
\label{sec:5.5}

下面，我们列出一些在机器学习中经常需要的常用梯度恒等式（Petersen and Pedersen, 2012）。这里，我们用 $ \mathrm{tr}(\cdot) $ 表示迹（参见定义 4.4），用 $ \det(\cdot) $ 表示行列式（参见第 4.1 节），并用 $ \boldsymbol{f}(\boldsymbol{X})^{-1} $ 表示 $ \boldsymbol{f}(\boldsymbol{X}) $ 的逆（假定其存在）。

\begin{equation}
\frac{\partial}{\partial \boldsymbol{X}} \boldsymbol{f}(\boldsymbol{X})^\top = \left( \frac{\partial \boldsymbol{f}(\boldsymbol{X})}{\partial \boldsymbol{X}} \right)^\top
\tag{5.99}\label{eq:5-99}
\end{equation}

\begin{equation}
\frac{\partial}{\partial \boldsymbol{X}} \mathrm{tr}(\boldsymbol{f}(\boldsymbol{X})) = \mathrm{tr}\left( \frac{\partial \boldsymbol{f}(\boldsymbol{X})}{\partial \boldsymbol{X}} \right)
\tag{5.100}\label{eq:5-100}
\end{equation}

\begin{equation}
\begin{adjustbox}{max width=\linewidth}
\frac{\partial}{\partial \boldsymbol{X}} \det(\boldsymbol{f}(\boldsymbol{X})) = \det(\boldsymbol{f}(\boldsymbol{X})) \, \mathrm{tr}\left( \boldsymbol{f}(\boldsymbol{X})^{-1} \frac{\partial \boldsymbol{f}(\boldsymbol{X})}{\partial \boldsymbol{X}} \right)
\end{adjustbox}
\tag{5.101}\label{eq:5-101}
\end{equation}

\begin{equation}
\begin{adjustbox}{max width=\linewidth}
\frac{\partial}{\partial \boldsymbol{X}} \boldsymbol{f}(\boldsymbol{X})^{-1} = -\boldsymbol{f}(\boldsymbol{X})^{-1} \frac{\partial \boldsymbol{f}(\boldsymbol{X})}{\partial \boldsymbol{X}} \boldsymbol{f}(\boldsymbol{X})^{-1}
\end{adjustbox}
\tag{5.102}\label{eq:5-102}
\end{equation}

\begin{equation}
\frac{\partial \boldsymbol{a}^\top \boldsymbol{X}^{-1} \boldsymbol{b}}{\partial \boldsymbol{X}} = -(\boldsymbol{X}^{-1})^\top \boldsymbol{a} \boldsymbol{b}^\top (\boldsymbol{X}^{-1})^\top
\tag{5.103}\label{eq:5-103}
\end{equation}

\begin{equation}
\frac{\partial \boldsymbol{x}^\top \boldsymbol{a}}{\partial \boldsymbol{x}} = \boldsymbol{a}^\top
\tag{5.104}\label{eq:5-104}
\end{equation}

\begin{equation}
\frac{\partial \boldsymbol{a}^\top \boldsymbol{x}}{\partial \boldsymbol{x}} = \boldsymbol{a}^\top
\tag{5.105}\label{eq:5-105}
\end{equation}

\begin{equation}
\frac{\partial \boldsymbol{a}^\top \boldsymbol{X} \boldsymbol{b}}{\partial \boldsymbol{X}} = \boldsymbol{a} \boldsymbol{b}^\top
\tag{5.106}\label{eq:5-106}
\end{equation}

\begin{equation}
\frac{\partial \boldsymbol{x}^\top \boldsymbol{B} \boldsymbol{x}}{\partial \boldsymbol{x}} = \boldsymbol{x}^\top (\boldsymbol{B} + \boldsymbol{B}^\top)
\tag{5.107}\label{eq:5-107}
\end{equation}

\begin{equation}
\begin{adjustbox}{max width=\linewidth}
\frac{\partial}{\partial \boldsymbol{s}} (\boldsymbol{x} - \boldsymbol{A} \boldsymbol{s})^\top \boldsymbol{W} (\boldsymbol{x} - \boldsymbol{A} \boldsymbol{s}) = -2(\boldsymbol{x} - \boldsymbol{A} \boldsymbol{s})^\top \boldsymbol{W} \boldsymbol{A} \quad (\textcn{对于对称矩阵 } \boldsymbol{W})
\end{adjustbox}
\tag{5.108}\label{eq:5-108}
\end{equation}

\emph{备注} 。在本书中，我们仅讨论矩阵的迹和转置。然而，我们已经看到导数可以是更高维的张量，在那种情况下通常的迹和转置并没有定义。在这些情况下，$ D \times D \times E \times F $ 维张量的迹将是一个 $ E \times F $ 维矩阵。这是 \emph{张量缩并（tensor contraction）} 的特例。类似地，当我们对张量进行“转置”时，我们的意思是对换前两个维度。具体而言，在式~\eqref{eq:5-99} 至式~\eqref{eq:5-102} 中，当我们处理多元函数 $ \boldsymbol{f}(\cdot) $ 并计算关于矩阵的导数时（并且选择不按照第 5.4 节讨论的那样将其向量化），需要进行与张量相关的计算。♦
